% Figura 3 de la postulación: balance de despeje y capacidad de la maquinaria y del personal.
\documentclass[border=4pt]{standalone}
% Preámbulo común de las figuras 1 a 4 de la postulación (se carga con \input en cada figura).
% Los cortes y umbrales vienen de umbrales_calibrados.tex (lo escribe figuras_postulacion.py).
\usepackage{fontspec}
\setmainfont{Cambria}
\usepackage[table]{xcolor}
\usepackage{array,booktabs}
\usepackage{tikz}
\usetikzlibrary{arrows.meta,positioning,calc,fit,backgrounds,decorations.pathreplacing}
\definecolor{fondo}{HTML}{F2F1EC}
\definecolor{txt2}{HTML}{52514E}
% estados, ordenados por magnitud: E3 despeje < E4 tormenta < E5 peligro
\definecolor{cE0}{HTML}{D9D8D2}\definecolor{cE1}{HTML}{1BAF7A}\definecolor{cE2}{HTML}{EDA100}
\definecolor{cE3}{HTML}{2A78D6}\definecolor{cE4}{HTML}{E34948}\definecolor{cE5}{HTML}{4A3AA7}
\definecolor{cSV}{HTML}{2A78D6}\definecolor{cH}{HTML}{EB6834}\definecolor{cN}{HTML}{52514E}
% Generado por figuras_postulacion.py desde salidas/sistema/parametros.json. No editar a mano.
\newcommand{\cortesS}{\textbf{0{,}5} & 5 & \textbf{12{,}5} & 20 & \textbf{25}}
\newcommand{\cortesF}{0{,}05 & 0{,}25 & 0{,}50 & 0{,}75 & 0{,}95}
\newcommand{\cortesr}{0{,}5 & 1 & 2 & 3 & 4}
\newcommand{\cortesdur}{1 & 6 & 12 & 18 & 23}
\newcommand{\cortesG}{10 & 15 & 20 & 25 & 30}
\newcommand{\cortesR}{\textbf{0{,}5} & 2{,}5 & 5 & \textbf{10} & 15}
\newcommand{\cortesHNS}{10 & 20 & 30 & 50 & 80}
\newcommand{\cortesRS}{10 & 20 & 40 & 60 & 100}
\newcommand{\pesoS}{0{,}45}
\newcommand{\pesoF}{0{,}15}
\newcommand{\pesor}{0{,}10}
\newcommand{\pesodur}{0{,}10}
\newcommand{\pesoG}{0{,}10}
\newcommand{\pesoR}{0{,}10}
\newcommand{\SVrest}{1}
\newcommand{\SVparc}{2{,}5}
\newcommand{\SVcierre}{3{,}5}
\newcommand{\Hrest}{3}
\newcommand{\Hparc}{4}
\newcommand{\Hcierre}{5}
\newcommand{\Nrest}{5}
\newcommand{\Nparc}{30}
\newcommand{\Ncierre}{50}
\newcommand{\cdesp}{6}
\newcommand{\expAlto}{1}
\newcommand{\expMedio}{0{,}5}
\newcommand{\hnSobrecarga}{30}

\newcolumntype{L}[1]{>{\raggedright\arraybackslash}p{#1}}
\newcolumntype{C}{>{\centering\arraybackslash}p{0.62cm}}
% encabezado de notas con degradado del color del índice
\newcommand{\notas}[1]{\cellcolor{#1!12}1 & \cellcolor{#1!28}2 & \cellcolor{#1!45}3 & \cellcolor{#1!65}\textcolor{white}{4} & \cellcolor{#1!85}\textcolor{white}{5}}
\tikzset{
  tit/.style={anchor=north west, font=\small\bfseries, inner sep=0pt},
  caja/.style={draw=#1, line width=0.9pt, rounded corners=3pt, align=center, minimum width=4.2cm, minimum height=1.05cm, fill=white},
  dato/.style={draw=txt2, dashed, rounded corners=3pt, align=center, minimum width=4.5cm, minimum height=0.85cm},
  est/.style={rounded corners=2pt, align=center, minimum width=1.98cm, minimum height=1.05cm, font=\scriptsize, text=white, inner sep=1pt},
  niv/.style={rounded corners=2pt, align=center, minimum width=2.45cm, minimum height=0.62cm, font=\scriptsize, inner sep=1pt},
  llave/.style={decorate, decoration={brace, amplitude=3pt}, draw=txt2},
  flecha/.style={-{Stealth[length=2mm]}, thick, draw=txt2},
  eti/.style={font=\scriptsize, color=txt2, align=center}}

\begin{document}
\begin{tikzpicture}[font=\footnotesize,
  paso/.style={draw=txt2, rounded corners=3pt, align=center, fill=white, font=\scriptsize, minimum height=0.9cm, inner sep=3pt}]
% a) Balance de despeje
\node[tit] at (0,0) {a) Balance de despeje, $N$ (cm de nieve en la calzada)};
\node[anchor=north west, inner sep=0pt] (td) at (0,-0.5) {\footnotesize\setlength{\tabcolsep}{3pt}\renewcommand{\arraystretch}{1.25}%
\begin{tabular}{>{\centering\arraybackslash}p{0.35cm}L{6.9cm}}
\toprule
& nieve que quedó en la calzada al cerrar el bloque anterior \\
$+$ & nieve que cae sobre el camino durante el bloque \\
$+$ & nieve que el viento deposita sobre el camino \\
$-$ & nieve que la maquinaria alcanza a sacar en 6 h, con al menos \cdesp\ cm/h (nada si $H = \Hcierre$, porque las cuadrillas no entran) \\
\midrule
$=$ & \textbf{$N$: nieve que queda por despejar} \\
\bottomrule
\end{tabular}};
% niveles, justo bajo la tabla
\path (td.south) ++(0,-0.45) coordinate (yc);
\coordinate (cx1) at (1.25,0); \coordinate (cx2) at (3.8,0); \coordinate (cx3) at (6.55,0);
\node[niv, fill=cE1, text=white] at (cx1 |- yc) {restricción\\$N > \Nrest$ cm};
\node[niv, fill=cE2] at (cx2 |- yc) {cierre parcial\\$N > \Nparc$ cm};
\node[niv, fill=cE3, text=white, minimum width=2.9cm] (d3) at (cx3 |- yc) {cerrado por despeje\\$N > \Ncierre$ cm (tras un cierre)};
\node[anchor=north west, align=left, font=\scriptsize, text width=7.9cm, inner sep=0pt] (ent) at ($(0,0 |- d3.south)+(0,-0.25)$)
  {\textbf{Para qué sirve:} la ruta puede reabrir cuando $N$ baja de \Ncierre\ cm. Con esa meta, el balance responde cada 6 h dos preguntas: cuántas horas faltan para reabrir (la espera hasta que baje el peligro más el tiempo que la maquinaria tarda en bajar $N$ a \Ncierre\ cm) y cuántos equipos y operadores hacen falta para reabrir en el plazo que fije la autoridad.};

% b) De la flota a la capacidad
\node[tit] at (8.55,0) {b) De la maquinaria y el personal a la capacidad};
\node[dato, minimum width=3.65cm, font=\scriptsize] (maq) at (10.5,-1.0) {\textbf{Maquinaria}: tipo, número\\y velocidad de avance en nieve};
\node[dato, minimum width=3.65cm, font=\scriptsize] (per) at (14.35,-1.0) {\textbf{Personal}: operadores\\por turno y horario};
\node[paso, minimum width=7.5cm] (ope) at (12.43,-2.45) {\textbf{Equipos operativos}: los que tienen operador en el turno};
\node[paso, minimum width=7.5cm, draw=cN, line width=0.9pt] (cap) at (12.43,-3.75) {\textbf{Capacidad de despeje} $c$ (cm/h) = equipos operativos $\times$ velocidad $\times$\\espesor que mantienen, dividido por pasadas $\times$ largo del tramo};
\node[paso, minimum width=3.65cm] (hrs) at (10.5,-5.35) {Horas para reabrir};
\node[paso, minimum width=3.65cm] (flo) at (14.35,-5.35) {Equipos y operadores para\\reabrir en un plazo dado};
\draw[flecha] (maq) -- (maq |- ope.north);
\draw[flecha] (per) -- (per |- ope.north);
\draw[flecha] (ope) -- (cap);
\draw[flecha] (cap.south -| hrs) -- (hrs);
\draw[flecha] (cap.south -| flo) -- (flo);
\draw[flecha] (cap.west) -- node[eti, above]{$c$} ($(td.east)+(0.05,-0.25)$);
\end{tikzpicture}
\end{document}
